% !TeX root = ../tesis.tex

Durante este capítulo, mantendremos la notación y la arquitectura fijadas en la Sección~\ref{sec:design-decisions-arquitectura}: $\ChainA$ denota la cadena de origen del evento a verificar y $\ChainB$ la cadena destino del \zkBridge. La implementación se presenta en dos bloques complementarios. Primero describimos la etapa que materializa la lógica \Onchain del camino $\LockTx \rightarrow \MintTx$ sobre el modelo \eUTxO usando validadores escritos en Aiken \cite{aiken2026, Chakravarty2020}. Luego incorporamos los circuitos \ZK concretos, la verificación de certificados de \Mithril y el flujo transaccional final. El camino inverso $\BurnTx \rightarrow \UnlockTx$ se incluye en la base de código, y en este trabajo se lo trata como una extensión natural para que la arquitectura del \zkBridge sea bidireccional, pero no lo consideramos como foco central de la implementación.

\section{Objetivo y alcance del bloque inicial}
\label{sec:validadores-y-flujo-transaccional-objetivo-y-alcance-del-bloque-inicial}

La Subsección~\ref{subsec:design-decisions-capas-actualizacion-y-aplicacion} estableció la separación entre una capa de actualización y una capa de aplicación. La primera mantiene en $\ChainB$ un estado autenticado y compacto sobre $\ChainA$, mientras que la segunda usa ese estado para autorizar operaciones concretas del \Bridge. La implementación de este capítulo baja esa división a componentes precisos: \UTxOs de estado identificados por NFTs, \textit{datums} que representan el estado persistente del protocolo, \textit{redeemers} que transportan la evidencia puntual de cada transacción, y validadores que fuerzan la atomicidad entre todos esos componentes.

La elección de Aiken como lenguaje de implementación resulta natural en este contexto. Por un lado, permite programar directamente scripts compatibles con el modelo \eUTxO. Por otro, ofrece una interfaz relativamente directa para integrar verificadores de \Groth sobre \BLS mediante bibliotecas como \texttt{ak-381}, lo cual encaja con la decisión arquitectónica de dejar el trabajo pesado \Offchain y reservar para la cadena destino una verificación sucinta \cite{CIP0381, IOGPlutusV3, AK381Repo}. En este bloque inicial, sin embargo, el objetivo no es cerrar todavía los circuitos de inclusión y de verificación de certificados, sino fijar correctamente las interfaces entre estado, argumentos y transacciones. Por eso aparecen \textit{placeholder circuits} en las posiciones donde el prototipo ya necesita un argumento \ZK; sus implementaciones completas se desarrollan en las Secciones~\ref{sec:tx-snapshot-membership-el-circuito-de-pertenencia-a-una-transactionsnapshot} y~\ref{sec:tx-set-update-el-circuito-txsetupdatecircuit}.

También es importante aclarar un punto metodológico. En una iteración intermedia del prototipo se experimentó con recomputar \Onchain el hash completo de certificados de \Mithril para encadenarlos dentro del script. Esa estrategia no se adopta como decisión final del sistema, ya que hashear certificados completos dentro del validador agrega demasiado \textit{overhead} en bytes procesados y \ExUnits, mientras que el \textit{binding} entre certificado y estado almacenado \Onchain puede lograrse de forma más barata y dentro de un circuito \Offchain. En consecuencia, en este capítulo describimos la implementación con la decisión final ya incorporada: el estado \Onchain publica los certificados y los encadena con la evidencia criptográfica necesaria, mientras que la validación detallada de esa evidencia se desplaza a circuitos especializados y a entradas públicas cuidadosamente elegidas.

\section{Componentes principales del prototipo}
\label{sec:validadores-y-flujo-transaccional-componentes-principales-del-prototipo}

La implementación se organiza alrededor de un conjunto pequeño de módulos y validadores, cada uno con una responsabilidad muy puntual. El repertorio central puede resumirse en la Tabla~\ref{tab:chapter6-validators}. La distinción más importante es entre validadores que gobiernan estado persistente y validadores que gobiernan una acción de aplicación. En la primera categoría quedan el validador de \StakeDistribution y el actualizador del conjunto de transacciones procesadas. En la segunda categoría queda, para esta etapa, el validador de minting.

\begin{table}[ht]
  \centering
  \scriptsize
  \begin{tabular}{p{3.5cm}p{2.0cm}p{8.0cm}}
    \toprule
    \textbf{Validador} & \textbf{Rol} & \textbf{Responsabilidad principal} \\
    \midrule
    \StakeDistributionMintValidator & Bootstrap & Mintear una única NFT de \StakeDistribution a partir de un certificado génesis y crear el \UTxO inicial de estado. \\
    \texttt{stake\_distribution.}\allowbreak\texttt{spend} & Actualización & Consumir el \UTxO que contiene el certificado padre y recrearlo con el nuevo certificado publicado por \StakeDistTx. \\
    \MintingValidator & Aplicación & Autorizar el \textit{minting} de tokens envueltos en $\ChainB$ demostrando que existe una $\LockTx$ válida en $\ChainA$ y que su evidencia no fue reutilizada. \\
    \TxsUpdaterMintingValidator & Anti-replay & Actualizar (en la misma transacción que el minting) la raíz del \SMT de transacciones ya procesadas. \\
    \UnlockingValidator y \TxsUpdaterUnlockingValidator & Extensión & Extender el mismo patrón al camino $\BurnTx \rightarrow \UnlockTx$, reservado aquí como versión bidireccional futura. \\
    \bottomrule
  \end{tabular}
  \caption{Validadores relevantes del bloque inicial de implementación.}
  \label{tab:chapter6-validators}
\end{table}

Alrededor de esos scripts aparecen tres módulos auxiliares particularmente importantes. El primero define los tipos de dominio compartidos por todo el sistema: el datum del \UTxO generado por la \LockTx, el datum mantenido por el \TxsUpdaterUnlockingValidator y los \textit{redeemers} que transportan identificadores de transacción, argumentos \ZK \Groth y referencias a \UTxOs. El segundo encapsula la verificación genérica de argumentos \ZK \Groth en Aiken. El tercero concentra funciones auxiliares para encontrar \UTxOs candidatos válidos, verificar la presencia de NFTs de estado y coordinar la coherencia entre redeemers de distintos scripts. El objetivo es evitar repetir lógica de \eUTxO y dejar aisladas las zonas donde luego se reemplazarán \textit{proof stubs} por verificadores reales.

\section{Representación del estado del bridge en \eUTxO}
\label{sec:validadores-y-flujo-transaccional-representacion-del-estado-del-bridge-en-eutxo}

El bloque inicial del prototipo concreta la decisión de la Sección~\ref{sec:design-decisions-estado-en-eutxo-atomicidad-y-concurrencia}: en lugar de modelar el estado del \zkBridge como un registro global mutable, lo representa mediante una pequeña familia de \UTxOs canónicos, cada uno protegido por un script y marcado por una NFT. Esta elección reaprovecha directamente las fortalezas del modelo \eUTxO \cite{Chakravarty2020}: cada transición de estado consume un \UTxO bien identificado y crea su sucesor, con lo cual la preservación de invariantes queda reducida a un chequeo local sobre la transacción que implementa el paso.

En el camino $\LockTx \rightarrow \MintTx$ aparecen tres piezas persistentes.

\begin{enumerate}[noitemsep]
  \item El \UTxO de fondos bloqueados en $\ChainA$. Su dirección de pago corresponde al script del \UnlockingValidator, y su datum contiene dos campos: el identificador del \zkBridge y la dirección donde deberán recibirse los \textit{wrapped assets} en $\ChainB$ (es decir, el beneficiario autorizado del \textit{minting}).
  \item El \UTxO del conjunto de transacciones ya procesadas en $\ChainB$. Su datum no almacena explícitamente una lista creciente de transacciones, sino la raíz de un \SMT. El NFT asociado identifica unívocamente este estado y obliga a que toda actualización consuma exactamente el \UTxO correcto.
  \item El \UTxO de \StakeDistribution en $\ChainB$. Su datum contiene el certificado de \Mithril más reciente aceptado por el protocolo. Al igual que en el caso anterior, un NFT identifica el estado vigente y fuerza a que cada actualización transfiera el control al mismo script.
\end{enumerate}

Esta descomposición responde a tres necesidades distintas. El \UTxO de fondos bloqueados los inmoviliza y deja registrada \Onchain la intención de transferencia; el \UTxO del conjunto de transacciones procesado resuelve el problema de \textit{anti-replay}; y el \UTxO de \StakeDistribution permite verificar que la evidencia presentada por un usuario en $\ChainB$ está conectada con una cadena válida de certificados de \Mithril. Es relevante observar que ninguna de estas estructuras necesita crecer linealmente con la historia del \zkBridge. El costo de mantener el estado se controla usando \textit{commitments} compactos y actualizaciones por sustitución, no por acumulación.

Además, la decisión tiene consecuencias positivas en términos de interfaz de transacción: cada paso del protocolo debe transportar solo la porción de información estrictamente necesaria para reconstruir el invariante local. Por eso los \textit{datums} contienen estado estable y reusable, mientras que los \textit{redeemers} contienen evidencia efímera: argumentos \ZK, identificadores de transacción y otros campos puntuales que permiten reconstruir la transacción bloqueante. Esto permite mantener los scripts pequeños y respetar \maxTxSize.

\section{Flujo transaccional del bloque inicial}
\label{sec:validadores-y-flujo-transaccional-flujo-transaccional-del-bloque-inicial}

La Figura~\ref{fig:chapter6-transaction-flow} resume el flujo que implementa el prototipo para el caso principal de uso. El orden lógico no es accidental: primero debe haber certificados autenticados de \StakeDistribution; luego puede ocurrir una $\LockTx$ en $\ChainA$; y recién entonces una transacción de \textit{minting} puede presentar la evidencia necesaria para emitir los \textit{wrapped assets} en $\ChainB$.

\begin{figure}[!t]
  \centering
  \shorthandoff{>}
  \begin{tikzpicture}[
    scale=0.84,
    transform shape,
    font=\small,
    box/.style={draw, rounded corners=4pt, thick, align=center, minimum width=3.1cm, minimum height=0.95cm},
    statebox/.style={draw=orange!70!black, fill=orange!12, rounded corners=4pt, thick, align=center, minimum width=3.4cm, minimum height=1cm},
    appbox/.style={draw=blue!70!black, fill=blue!10, rounded corners=4pt, thick, align=center, minimum width=3.3cm, minimum height=1cm},
    offbox/.style={draw=gray!70!black, fill=gray!12, rounded corners=4pt, thick, align=center, minimum width=3.2cm, minimum height=0.95cm},
    lbl/.style={fill=white, inner sep=1pt}
  ]
    \node at (-6, 3.2) {\textbf{Cadena origen $\ChainA$}};
    \node at (4, 3.2) {\textbf{Cadena destino $\ChainB$}};

    \node[appbox] (lock) at (-6, 1.7) {$\LockTx$\\bloquea activos};
    \node[statebox] (locked) at (-6, 0.0) {\UTxO bloqueado\\\UnlockingValidator};

    \node[offbox] (relayer) at (-3.0, -2.5) {\textit{relayer/prover}\\recolecta evidencia};

    \node[appbox] (sdtx) at (4, 2.0) {$\StakeDistTx$\\publica nuevo certificado};
    \node[statebox] (sdstate) at (4, 0) {\UTxO de \StakeDistribution\\+ NFT de estado};
    \node[statebox] (txset) at (4, -2) {\UTxO del conjunto usado\\raíz \SMT + NFT};
    \node[appbox] (mint) at (4, -4) {$\MintTx$\\mintea tokens envueltos};

    \draw[->, thick] (lock) -- node[lbl, right] {crea} (locked);
    \draw[->, thick] (locked.south) -- node[lbl, above] {tx id, datum, valor} (relayer.west);
    \draw[->, thick] (relayer.east) -- node[lbl, above] {certificado + argumento \ZK} (sdstate.west);
    \draw[->, thick] (sdtx) -- node[lbl, right] {actualiza} (sdstate);
    \draw[->, thick] (sdstate.east) .. controls (7.2,0.6) and (7.2,-1.7) .. node[lbl] {\UTxO \StakeDistribution (ref)} (mint.east);
    
    \draw[->, thick] (txset.south) -- node[lbl, right] {NFT} (mint.north);
    \draw[->, thick] (relayer.south east) -- node[lbl, below left] {argumento \ZK de inclusión} (mint.west);
    \draw[->, thick] (mint) -- ++(0,-1.2) node[lbl, below] {\UTxO al destinatario};
    \draw[->, thick] (mint.north west) -- node[lbl] {actualiza \SMT} (txset.south west);
  \end{tikzpicture}
  \shorthandon{>}
  \caption{Flujo transaccional del bloque inicial del \zkBridge para el camino $\LockTx \rightarrow \MintTx$. La \LockTx crea un \UTxO bloqueado en la cadena origen; el \textit{relayer}/\textit{prover} recolecta el identificador, datum, valor y evidencia criptográfica; en la cadena destino, \StakeDistTx actualiza el \UTxO de \StakeDistribution y \MintTx mintea tokens envueltos mientras actualiza el conjunto anti-replay resumido por una raíz \SMT.}
  \label{fig:chapter6-transaction-flow}
\end{figure}

\subsection{Actualización de \StakeDistribution mediante \texorpdfstring{\StakeDistTx}{tx\_sd}}
\label{subsec:validadores-y-flujo-transaccional-actualizacion-de-stakedistribution-mediante-tx-sd}

La transacción \StakeDistTx pertenece a la capa de actualización. Su objetivo no es mover activos del usuario, sino mantener en $\ChainB$ un resumen autenticado del estado de \Mithril que luego será utilizado por las transacciones de minting. Conceptualmente, el validador admite dos fases.

En la fase de \textit{bootstrap}, un script de minting crea el NFT que identificará al estado de \StakeDistribution y lo deposita en un \UTxO cuyo datum contiene un certificado génesis de \Mithril. Esto debe poder ocurrir una sola vez; por eso se la ata a un único \UTxO capaz de acuñar el NFT, y se exige que el \UTxO creado quede bajo el control del mismo validador, para que gobierne las actualizaciones posteriores.

En la fase recurrente de actualización, \StakeDistTx consume el \UTxO que contiene el certificado padre y produce un nuevo \UTxO con el mismo NFT y con el certificado nuevo en el datum. Este bloque inicial ya fija varios invariantes útiles: el nuevo certificado debe ser del tipo correcto, debe estar correctamente encadenado con su padre y el control del estado debe permanecer en el mismo script. La interpretación criptográfica de una actualización correcta se desarrolla en la Sección~\ref{sec:mithril-cadena-de-certificados}; la Sección~\ref{sec:mithril-a-cardano-del-circuito-halotwo-de-mithril-a-cardano} completa cómo esa evidencia se consume mediante recibos de verificación.

Un detalle operativo importante es que el \UTxO de \StakeDistribution no se consume en cada \MintTx. En lugar de eso, aparece como \textit{reference input}. Esta decisión reduce costos y evita que cada usuario compita por modificar un estado que solamente necesita leer.

\subsection{Transacción de \textit{locking} en la cadena de origen}
\label{subsec:validadores-y-flujo-transaccional-transaccion-de-locking-en-la-cadena-de-origen}

La $\LockTx$ materializa la intención del usuario de iniciar un traslado de valor. En este bloque inicial, la transacción crea un \UTxO dirigido al script del \UnlockingValidator de $\ChainA$, el cual contiene la totalidad de los activos bloqueados y un datum con dos componentes: \texttt{bridge\_id} y \texttt{destination\_address}. El primero evita que la misma transacción sea reutilizada para reclamar activos en múltiples \Bridges incompatibles, mientras que el segundo deja fijado desde $\ChainA$ quién es el destinatario autorizado al hacer \MintTx en $\ChainB$.

El campo \texttt{destination\_address} debe leerse como una dirección de la cadena destino, codificada dentro del datum de la transacción de origen. No asumimos que el espacio de direcciones de $\ChainA$ coincida con el de $\ChainB$, ni que $\ChainA$ pueda interpretar semánticamente esa dirección: para $\ChainA$ es simplemente un dato transportado por el datum, mientras que su significado se recupera al validar \MintTx en $\ChainB$.

La implementación adopta en este bloque inicial una convención deliberadamente conservadora: el \UTxO bloqueado se toma como la parte semántica relevante de la operación, y la evidencia presentada en $\MintTx$ se organiza alrededor de ese \UTxO. La información útil para validación queda concentrada en estos campos: referencia del input, dirección de la entrada, datum original, dirección de destino en $\ChainB$ y valor transferido. Esto permite expresar con claridad qué propiedades debe preservar el \zkBridge entre la transacción en $\ChainA$ remota y el minting en $\ChainB$, incluso antes de fijar la forma definitiva del circuito de inclusión.

En esta etapa, el prototipo trabaja sobre una formulación isomorfa del problema, donde las cadenas $\ChainA, \ChainB$ son isomorfas. Esto significa que la interfaz de las transacciones en ambas cadenas puede representarse en Aiken con estructuras suficientemente similares como para chequear localmente la coherencia de la evidencia presentada sobre la \LockTx en el script de \MintingValidator. Entonces, la capa de aplicación del \zkBridge puede validarse antes de introducir las complejidades propias de un circuito completo de inclusión de transacciones de \Cardano. En la Sección~\ref{sec:tx-snapshot-membership-el-circuito-de-pertenencia-a-una-transactionsnapshot}, esta responsabilidad se separa: el circuito fija criptográficamente que el \TxId de la \LockTx pertenece a \CardanoTransactions, pero no prueba por sí solo el destinatario ni el valor transferido. Por eso, la coherencia semántica entre el evento de origen y el minting queda como invariante explícito de la interfaz del \Bridge y de los validadores que consumen esa evidencia.

\subsection{Transacción de minting en la cadena destino}
\label{subsec:validadores-y-flujo-transaccional-transaccion-de-minting-en-la-cadena-destino}

La \MintTx es la pieza central de la capa de aplicación. Su trabajo consiste en \textit{mintear} en $\ChainB$ un valor equivalente al bloqueado en $\ChainA$, pero solo si la evidencia aportada por el usuario satisface simultáneamente cuatro condiciones:

\begin{enumerate}[noitemsep]
  \item La \LockTx efectivamente pertenece a $\ChainA$ y fue incluida en un \TransactionSnapshot autenticado por \Mithril.
  \item El hash o identificador de la \LockTx no fue usado antes para mintear.
  \item El valor \textit{minteado} del \textit{wrapped asset} coincide con el valor bloqueado en origen.
  \item La dirección de destino en $\ChainB$ codificada en \LockTx coincide con el receptor efectivo de los \textit{wrapped assets}.
\end{enumerate}

Para verificar esto, el \textit{redeemer} de \MintTx asociado a \MintingValidator transporta dos componentes esenciales. La primera es la evidencia ligada a \Mithril: un certificado de \TransactionSnapshot, un \ProofReceipt del dominio \texttt{cardano\_transactions} que autentica la verificación \STM de ese certificado, y un argumento \Groth de inclusión del \TxId de \LockTx dentro del \TransactionSnapshot. La segunda es la información de aplicación necesaria para ligar ese identificador con el reclamo concreto que se está minteando.

La implementación del \MintingValidator sigue entonces una estrategia de doble chequeo. Por un lado, verifica la coherencia externa de la evidencia: existe un certificado aceptable, el \TransactionSnapshot contiene la transacción indicada y el conjunto anti-replay será actualizado en la misma operación. Por otro lado, verifica la coherencia interna: con los campos provistos en el \textit{redeemer} y con el valor realmente \textit{minteado} por la transacción, reconstruye una versión esperada de la \LockTx y comprueba que su identificador coincida con el reclamado. No alcanza con que exista alguna transacción en el \TransactionSnapshot, sino que debe tratarse exactamente de \LockTx.

Esta última condición se complementa con el mecanismo \textit{anti-replay}. Una vez que existe un camino para probar inclusión, el siguiente problema natural es evitar el reprocesamiento de la misma evidencia. La implementación lo resuelve mediante un \SMT cuya raíz se almacena en el datum de un \UTxO de estado. Cada hoja del árbol representa un \textit{commitment} a una transacción ya utilizada por el \zkBridge. La decisión de usar un \SMT en lugar de una lista explícita responde directamente al crecimiento no acotado del historial: el conjunto puede expandirse indefinidamente sin volver inviable el almacenamiento \Onchain. La actualización de ese estado está a cargo del validador \TxsUpdaterMintingValidator. Su \textit{redeemer} incluye el identificador de la \LockTx y del \UTxO generado por ella, así como un argumento \ZK \Groth (simulado en este capítulo) que certifica la correcta actualización de la raíz del \SMT con la inclusión de la nueva transacción.

Lo importante es que \MintingValidator y \TxsUpdaterMintingValidator no se validan en transacciones separadas, sino dentro de la misma \MintTx. El protocolo no permite que una transacción agregue un identificador al conjunto usado sin mintear, ni que haga un \textit{minting} sin actualizar ese conjunto. En el modelo \eUTxO esto se logra obligando a que ambos scripts participen de la misma operación y comprobando la presencia del otro desde cada uno de ellos. Es precisamente este tipo de decisión la que vuelve defendible al diseño en el marco del modelo \eUTxO.

\begin{itemize}[noitemsep]
	\item Desde \MintingValidator, el script exige que entre los inputs exista el \UTxO que porta el NFT asociado a \TxsUpdaterMintingValidator.
	\item Desde \TxsUpdaterMintingValidator, el script valida que la transacción está efectivamente minteando tokens del \zkBridge.
	\item Ambos scripts comparan el contenido de sus respectivos \textit{redeemers}, de modo de comprobar que se están refiriendo a la misma \LockTx.
\end{itemize}

\section{Decisiones de ingeniería y preparación para los circuitos reales}
\label{sec:validadores-y-flujo-transaccional-decisiones-de-ingenieria-y-preparacion-para-los-circuitos-reales}

El bloque inicial no solo implementa validadores; también fija varias decisiones prácticas que condicionan la evolución del sistema. La primera es el uso sistemático de NFTs para marcar estado canónico. Tanto el conjunto de transacciones usadas como el estado de \StakeDistribution dependen de una NFT que identifica de manera unívoca al \UTxO válido. Esta técnica evita ambigüedades sobre qué \textit{output} representa el estado vigente y simplifica mucho la lógica de los scripts.

La segunda es la separación entre datos persistentes y evidencia efímera. Los certificados de \Mithril y las raíces comprometidas viven en datums de \UTxOs reutilizables; los argumentos \ZK, identificadores de transacción y reconstrucciones parciales de la \LockTx viven en redeemers. Esa división es saludable para el costo del protocolo: aquello que muchos usuarios deben leer se publica una sola vez; aquello que cambia en cada operación viaja solo cuando la operación lo necesita.

La tercera, ya mencionada, es la decisión final de no recomputar \Onchain hashes completos de certificados de \Mithril. Desde el punto de vista teórico, tal recomputación podría parecer una forma directa de lograr \textit{binding} entre el dato publicado y el argumento \ZK presentado. Sin embargo, en la práctica introduce un costo desproporcionado en serialización, hashing y \ExUnits. La versión final del sistema persigue el mismo objetivo de un modo más económico: el \UTxO de \StakeDistribution publica el certificado relevante, los argumentos quedan ligados a entradas públicas del circuito que identifican el estado certificado y los validadores solo verifican la consistencia necesaria para que ese enlace sea unívoco. El desarrollo criptográfico de este punto se completa en las Secciones~\ref{sec:mithril-a-cardano-del-circuito-halotwo-de-mithril-a-cardano}, \ref{sec:tx-snapshot-membership-el-circuito-de-pertenencia-a-una-transactionsnapshot} y~\ref{sec:tx-set-update-el-circuito-txsetupdatecircuit}, pero lo mencionamos aquí porque condiciona la forma final de la interfaz \Onchain.

Una cuarta decisión es el orden de compilación e instanciación de scripts. Dado que varios validadores necesitan conocer \textit{policy IDs}, \textit{script hashes} o referencias únicas producidas por otros componentes, el prototipo induce un pequeño grafo de dependencias entre validadores. Entonces, este orden es consecuencia directa de cómo los scripts se referencian mutuamente. En la parte del sistema correspondiente a $\LockTx \rightarrow \MintTx$, dicho grafo puede resumirse como:

\begin{enumerate}[noitemsep]
  \item \StakeDistributionMintValidator y la política de minting de la NFT del conjunto de transacciones usadas se compilan primero, porque fijan los \textit{policy IDs} de los dos \UTxOs de estado sobre los que descansa el resto del protocolo.
  \item \StakeDistributionSpendValidator se compila en un segundo paso, ya que depende de la NFT de \StakeDistribution definida en el paso anterior para reconocer el \UTxO padre y su sucesor.
  \item \MintingValidator también se compila en ese segundo paso, porque necesita conocer tanto la política asociada al estado de \StakeDistribution como la política de la NFT del actualizador del conjunto usado.
  \item \TxsUpdaterMintingValidator se compila al final, ya que depende de la política del actualizador, pero también del identificador del propio \MintingValidator, al que debe quedar enlazado para forzar la atomicidad de \MintTx.
\end{enumerate}

Por último, la base de código ya incorpora un \textit{test suite} en Aiken para cada uno de los validadores utilizados. Desde el punto de vista de la tesis, estos tests no reemplazan la evaluación de resultados, pero sí aportan una primera validación estructural sobre los casos positivos y negativos donde se rechazan certificados mal encadenados, argumentos \ZK inválidos o transacciones que intentan mintear sin actualizar el conjunto \textit{anti-replay}.

\section{Formalización de los validadores principales}
\label{sec:validadores-y-flujo-transaccional-formalizacion-de-los-validadores-principales}

En esta sección fijamos una notación más abstracta para los validadores principales del camino $\LockTx \rightarrow \MintTx$. Seguimos la distinción introducida al presentar el modelo \eUTxO entre \textit{spending validators} y \textit{minting validators}. Para los primeros utilizamos la forma \SpendValidatorSchema, mientras que para los segundos utilizamos la forma \MintValidatorSchema. En todos los casos, la igualdad a \True debe leerse como ``el validador acepta la transición codificada por la transacción''.

Dos aclaraciones ayudan a leer correctamente esta formalización:

\begin{itemize}[noitemsep]
  \item La notación omite parámetros auxiliares del código Aiken, como referencias locales al input gobernado o constantes de entorno, cuando estos pueden recuperarse del contexto de transacción o forman parte de la instanciación del script.
  \item El foco está puesto en la semántica de aceptación del validador, no en la organización exacta del código ni en funciones auxiliares de búsqueda sobre `inputs`, `outputs`, `reference inputs` o `redeemers`.
\end{itemize}

En esta formalización nos concentramos en cuatro validadores: \StakeDistributionMintValidator, \StakeDistributionSpendValidator, \MintingValidator y \TxsUpdaterMintingValidator. Los validadores del camino inverso pueden formalizarse de manera análoga. 

\subsection{Validador de bootstrap de \StakeDistribution}
\label{subsec:validadores-y-flujo-transaccional-validador-de-bootstrap-de-stakedistribution}

\begin{center}
\MintValidatorNotation{\StakeDistributionMintValidator}
\end{center}

La expresión anterior vale si y sólo si se satisfacen simultáneamente las siguientes condiciones:

\begin{enumerate}[noitemsep]
  \item \tx consume el único \UTxO $o$ autorizado para bootstrap del NFT de \StakeDistribution.
  \item \redeemer contiene un certificado génesis \GenesisCertificate de \Mithril.
  \item \tx crea un \UTxO $o^*$ cuyo datum $\datum^*$ contiene \GenesisCertificate.
  \item $o^*$ porta el NFT de \StakeDistribution bajo el \textsf{PolicyId} del \MintingValidator.
  \item \GenesisCertificate es válido según la clave de verificación génesis \GenesisVerificationKey fijada por el protocolo.
\end{enumerate}

\subsection{Validador de actualización de \StakeDistribution}
\label{subsec:validadores-y-flujo-transaccional-validador-de-actualizacion-de-stakedistribution}

\begin{center}
\SpendValidatorNotation{\StakeDistributionSpendValidator}
\end{center}

La expresión anterior vale si y sólo si se satisfacen simultáneamente las siguientes condiciones:

\begin{enumerate}[noitemsep]
  \item \datum contiene un certificado padre \StakeDistributionCertificate de \Mithril.
  \item \redeemer contiene un nuevo certificado \StakeDistributionCertificateNew y un argumento \ZK asociado a su validación.
  \item \tx crea un \textit{output} sucesor $o^*$ cuyo datum  $\datum^*$ contiene exactamente \StakeDistributionCertificateNew.
  \item $o^*$ queda bloqueado por el mismo script (el \StakeDistributionSpendValidator).
  \item $o$ y $o^*$ portan el mismo NFT de \StakeDistribution.
  \item \StakeDistributionCertificateNew es efectivamente un certificado de \StakeDistribution de \Mithril y está correctamente encadenado con su padre \StakeDistributionCertificate.
  \item El argumento \ZK asociado a la validación de \StakeDistributionCertificateNew es aceptado.
\end{enumerate}

\subsection{Validador de \textit{minting} del \zkBridge}
\label{subsec:validadores-y-flujo-transaccional-validador-de-minting-del-zkbridge}

\begin{center}
\MintValidatorNotation{\MintingValidator}
\end{center}

La expresión anterior vale si y sólo si se satisfacen simultáneamente las siguientes condiciones:

\begin{enumerate}[noitemsep]
  \item \tx incluye como \textit{reference input} el \UTxO canónico de \StakeDistribution \StakeDistributionUTxO, del cual se extrae el certificado padre \StakeDistributionCertificate necesario para validar el \TransactionSnapshot.
  \item \redeemer contiene un certificado \TransactionSnapshotCertificate de \TransactionSnapshot, un \ProofReceipt del dominio \texttt{cardano\_transactions} y un argumento \Groth de inclusión del \TxId de \LockTx dentro de \TransactionSnapshotCertificate.
  \item El \ProofReceipt autentica que \TransactionSnapshotCertificate fue aceptado por la verificación \STM de fase 2 y que su \texttt{statement\_hash} coincide con el mensaje firmado del certificado.
  \item El argumento \Groth de inclusión prueba que el identificador reclamado de \LockTx pertenece al \TransactionSnapshotCertificate autenticado.
  \item \tx consume el \UTxO identificado por el NFT del actualizador del conjunto de transacciones procesadas.
  \item El identificador remoto procesado por el actualizador coincide con el identificador reclamado por el \MintingValidator.
  \item En la formulación isomorfa de este bloque inicial, a partir de \redeemer y del valor \textit{minteado} por la transacción puede reconstruirse y verificarse una versión esperada de \LockTx. En la versión final, esta condición debe leerse como un invariante de interfaz: la prueba de pertenencia fija el \TxId, no todos los campos de aplicación.
  \item El valor minteado y la dirección de destino en $\ChainB$ coinciden con el valor bloqueado y el destinatario autorizado por \LockTx, cuando esos datos forman parte del reclamo de aplicación presentado al validador.
\end{enumerate}

\subsection{Validador del conjunto de transacciones procesadas}
\label{subsec:validadores-y-flujo-transaccional-validador-del-conjunto-de-transacciones-procesadas}

\begin{center}
\SpendValidatorNotation{\TxsUpdaterMintingValidator}
\end{center}

La expresión anterior vale si y sólo si se satisfacen simultáneamente las siguientes condiciones:

\begin{enumerate}[noitemsep]
  \item \datum contiene la raíz anterior del \SMT $\MerkleRoot{\MerkleTree{T}}$ de transacciones procesadas.
  \item \tx crea un \textit{output} $o^*$ cuyo datum $\datum^*$ contiene la nueva raíz $\MerkleRoot{\MerkleTree{T}}^*$ del \SMT, de modo que queda asociado al conjunto de transacciones procesadas.
  \item $o^*$ queda bloqueado por el mismo script (el \TxsUpdaterMintingValidator).
  \item \redeemer contiene el identificador de \LockTx, la referencia al \UTxO generado por ella y un argumento \ZK de actualización del \SMT.
  \item El argumento \ZK de actualización certifica que la nueva raíz resulta de insertar correctamente la transacción reclamada en el conjunto anterior.
  \item \tx está efectivamente minteando tokens del \zkBridge.
  \item Los identificadores de \LockTx de \redeemer y del \textit{redeemer} $\redeemer^*$ asociado al \MintingValidator coinciden (porque en \MintTx se verifican ambos validadores).
\end{enumerate}
